\begin{table}[htbp]
\centering
\caption{Robustness Checks: Sensitivity of Key Results to Alternative Assumptions}
\label{tab:robustness}
\footnotesize
\begin{tabular}{lrrrr}
\toprule
Assumption & SV 2\textdegree{}C & SV 1.5\textdegree{}C & Bank Loss & GDP Impact \\
 & (\$B) & (\$B) & (\$M) & (\%) \\
\midrule
  \textbf{Baseline} & \textbf{79.8} & \textbf{128.3} & \textbf{3,147} & \textbf{0.54\%} \\
  +2pp discount rate & 67.0 & 108.1 & 3,147 & 0.54\% \\
  $-$2pp discount rate & 96.1 & 154.2 & 3,147 & 0.54\% \\
  +10\% production volume & 87.7 & 141.2 & 3,147 & 0.54\% \\
  $-$10\% production volume & 71.8 & 115.5 & 3,147 & 0.54\% \\
  Coal price path +15\% & 91.7 & 147.6 & 3,147 & 0.54\% \\
  Coal price path $-$15\% & 67.8 & 109.1 & 3,147 & 0.54\% \\
  Alternative LGD (0.55) & 79.8 & 128.3 & 3,316 & 0.54\% \\
  Alternative credit mult.~(12$\times$) & 79.8 & 128.3 & 3,147 & 0.70\% \\
\bottomrule
\end{tabular}

\smallskip
\noindent \textit{Notes:} SV = aggregate stranded value in USD billions. Bank Loss = total banking sector losses in USD millions under the 1.5\textdegree{}C scenario. GDP Impact = total GDP impact under 1.5\textdegree{}C. Baseline uses a 10\% WACC, LGD = 0.45, and credit multiplier of 8$\times$. Each row varies one assumption while holding others at baseline. Discount rate sensitivity captures the effect on present-value calculations. LGD = loss given default applied to distressed coal exposures. Credit multiplier scales the GDP impact of bank deleveraging.
\end{table}
